\begin{abstract}
% Multi-teacher on-policy distillation combines specialized teachers in a shared student, but its learning dynamics depend on how training examples and teacher feedback are weighted. We study three questions: how loss averaging affects capability balance, how concentrated parameter changes are under joint training, and how vocabulary supervision affects those changes and the resulting capabilities. We first separate two effects of response length: weighting domains by their token counts and weighting responses within each domain. Controlled comparisons show that removing these effects leads to different capability tradeoffs. Joint training produces concentrated cumulative changes in stored model weights, while the corresponding full-precision optimizer weights change more broadly. At a common student and optimizer state, sampled-token, top-$k$, and full-vocabulary supervision produce similar update concentration. Along online training trajectories, sampled-token and top-$k$ supervision yield different domain scores. Together, these results connect the averaging rule to capability allocation and distinguish vocabulary coverage from concentration of parameter changes.

Multi-teacher on-policy distillation (MOPD) provides a general framework for consolidating the capabilities of multiple teachers specialized through reinforcement learning into a single student model. However, it remains unclear how teacher supervision is weighted across tasks, how it shapes the resulting optimization dynamics, and whether different teachers induce overlapping parameter updates. We study how loss averaging and supervision density shape its optimization dynamics and capability transfer. Our analysis identifies two distinct effects of response length under token averaging: changing the relative weights of domains and weighting responses unequally within each domain. Fixing domain weights removes only the first effect, while equal response weighting also removes the second. Controlled comparisons show that these choices change gradient direction and redistribute capability gains, with global-token averaging favoring mathematics and equal response weighting favoring instruction following in our experiments. We further find that joint training produces concentrated cumulative changes in FP16 model weights, while the FP32 weights change more broadly. Local comparisons at a common student and optimizer state reveal similar update concentration across sampled-token, top-k, and full-vocabulary supervision, despite their different vocabulary coverage. Online sampled-token and top-k training nevertheless develop different capability profiles across domains and checkpoints. These findings characterize how loss averaging allocates teacher supervision and show why vocabulary coverage alone is insufficient to characterize parameter-update concentration or capability transfer.
\end{abstract}
